% GENERATED FILE. Edit canonical lessons, then run npm run generate:books.
% canonical-source-hash: fnv1a64:06e228e540681e2c
% canonical-lessons: TE-C387-vimarsa, TE-C387-vimarsincu, TE-C387-prasamsincu, TE-C387-vyakhyanincu, TE-C387-tappubattu

\chapter{Praise and Criticism}
\label{ch:te-praise-and-criticism}
\hlchaptermodality{387}

\begin{chapteropening}
\textbf{By the end of this chapter:} \emph{I can praise, criticise and find fault.}

Complete the last lesson of chapter 387: I can praise, criticise and find fault.
\end{chapteropening}

\section[\texorpdfstring{\te{విమర్శ} (vimarśa)}{విమర్శ (vimarśa)}]{\te{విమర్శ} (vimarśa) — criticism; a critique}

\label{lesson:TE-C387-vimarsa}



\begin{quote}
Before the new one: say the Telugu for opposition, then the Telugu for to reject.
\end{quote}

\subsection*{You'll want to know: \te{విమర్శ}}
\textbf{\te{విమర్శ}} — \emph{vimarśa} — \textquotedblleft{}criticism; a critique\textquotedblright{}.

\begin{quote}
\te{మంచి} \te{విమర్శ} \te{మనకు} \te{సహాయం} \te{చేస్తుంది}.
\end{quote}

\begin{quote}
\emph{mañci vimarśa manaku sahāyaṁ cēstundi.}
\end{quote}

\begin{quote}
Good criticism helps us.
\end{quote}

Plural: \te{విమర్శలు}, \emph{vimarśalu}.

Not only blame: a review of a book or a film is also a \te{విమర్శ}.

\begin{cousinweb}
\textbf{\te{విమర్శ}} (\emph{vimarśa}) is borrowed from Sanskrit \emph{vimarśa}, \textquotedblleft{}consideration, deliberation, critical test\textquotedblright{} (Monier-Williams).
\end{cousinweb}

\begin{grammarlens}[title={What you've built}]
One more word for giving reasons and opinions.
\end{grammarlens}

\subsection*{Your turn}
\begin{itemize}
\raggedright
  \item \emph{Say it:} \emph{vimarśa}
  \item \emph{Say it:} \emph{vimarśa}, once more
  \item \emph{Say it:} \emph{mañci vimarśa manaku sahāyaṁ cēstundi.} — Good criticism helps us.
  \item \emph{Recall:} say the Telugu for opposition, then the Telugu for to reject, then say \emph{vimarśa} again
\end{itemize}

\begin{tcolorbox}[breakable,colback=teal!4,colframe=teal!35!black,title={Before you move on}]
What does \te{విమర్శ} mean? (\textquotedblleft{}Criticism; a critique\textquotedblright{}.) Say it once more.
\end{tcolorbox}
\section[\texorpdfstring{\te{విమర్శించు} (vimarśiñcu)}{విమర్శించు (vimarśiñcu)}]{\te{విమర్శించు} (vimarśiñcu) — to criticise}

\label{lesson:TE-C387-vimarsincu}



\begin{quote}
Before the new one: say the Telugu for to reject, then the Telugu for criticism.
\end{quote}

\subsection*{You'll want to know: \te{విమర్శించు}}
\textbf{\te{విమర్శించు}} — \emph{vimarśiñcu} — \textquotedblleft{}to criticise\textquotedblright{}.

\begin{quote}
\te{పత్రికలు} \te{ప్రభుత్వాన్ని} \te{విమర్శించాయి}.
\end{quote}

\begin{quote}
\emph{patrikalu prabhutvānni vimarśiñcāyi.}
\end{quote}

\begin{quote}
The papers criticised the government.
\end{quote}

Past (I …): \te{విమర్శించాను}, \emph{vimarśiñcānu}.

The person or thing criticised takes -\te{ని}: \te{ప్రభుత్వాన్ని} \te{విమర్శించు}.

\begin{cousinweb}
\textbf{\te{విమర్శించు}} (\emph{vimarśiñcu}) is built on Sanskrit \emph{vimarśa}, \textquotedblleft{}consideration, critical test\textquotedblright{} (Monier-Williams), with the Telugu verb ending -\te{ఇంచు}; the noun is \te{విమర్శ}.
\end{cousinweb}

\begin{grammarlens}[title={What you've built}]
One more word for giving reasons and opinions.
\end{grammarlens}

\subsection*{Your turn}
\begin{itemize}
\raggedright
  \item \emph{Say it:} \emph{vimarśiñcu}
  \item \emph{Say it:} \emph{vimarśiñcu}, once more
  \item \emph{Say it:} \emph{patrikalu prabhutvānni vimarśiñcāyi.} — The papers criticised the government.
  \item \emph{Recall:} say the Telugu for to reject, then the Telugu for criticism, then say \emph{vimarśiñcu} again
\end{itemize}

\begin{tcolorbox}[breakable,colback=teal!4,colframe=teal!35!black,title={Before you move on}]
What does \te{విమర్శించు} mean? (\textquotedblleft{}To criticise\textquotedblright{}.) Say it once more.
\end{tcolorbox}
\section[\texorpdfstring{\te{ప్రశంసించు} (praśaṁsiñcu)}{ప్రశంసించు (praśaṁsiñcu)}]{\te{ప్రశంసించు} (praśaṁsiñcu) — to laud; to praise}

\label{lesson:TE-C387-prasamsincu}



\begin{quote}
Before the new one: say the Telugu for criticism, then the Telugu for to criticise.
\end{quote}

\subsection*{You'll want to know: \te{ప్రశంసించు}}
\textbf{\te{ప్రశంసించు}} — \emph{praśaṁsiñcu} — \textquotedblleft{}to laud; to praise\textquotedblright{}.

\begin{quote}
\te{అందరూ} \te{ఆమె} \te{ధైర్యాన్ని} \te{ప్రశంసించారు}.
\end{quote}

\begin{quote}
\emph{andarū āme dhairyānni praśaṁsiñcāru.}
\end{quote}

\begin{quote}
Everyone praised her courage.
\end{quote}

Past (I …): \te{ప్రశంసించాను}, \emph{praśaṁsiñcānu}.

More formal than \te{పొగడు}: for speeches, ceremonies and the papers.

\begin{cousinweb}
\textbf{\te{ప్రశంసించు}} (\emph{praśaṁsiñcu}) is built on Sanskrit \emph{praśaṁsā}, \textquotedblleft{}praise, commendation\textquotedblright{} (Monier-Williams), with the Telugu verb ending -\te{ఇంచు}.
\end{cousinweb}

\begin{grammarlens}[title={What you've built}]
One more word for giving reasons and opinions.
\end{grammarlens}

\subsection*{Your turn}
\begin{itemize}
\raggedright
  \item \emph{Say it:} \emph{praśaṁsiñcu}
  \item \emph{Say it:} \emph{praśaṁsiñcu}, once more
  \item \emph{Say it:} \emph{andarū āme dhairyānni praśaṁsiñcāru.} — Everyone praised her courage.
  \item \emph{Recall:} say the Telugu for criticism, then the Telugu for to criticise, then say \emph{praśaṁsiñcu} again
\end{itemize}

\begin{tcolorbox}[breakable,colback=teal!4,colframe=teal!35!black,title={Before you move on}]
What does \te{ప్రశంసించు} mean? (\textquotedblleft{}To laud; to praise\textquotedblright{}.) Say it once more.
\end{tcolorbox}
\section[\texorpdfstring{\te{వ్యాఖ్యానించు} (vyākhyāniñcu)}{వ్యాఖ్యానించు (vyākhyāniñcu)}]{\te{వ్యాఖ్యానించు} (vyākhyāniñcu) — to comment; to remark}

\label{lesson:TE-C387-vyakhyanincu}



\begin{quote}
Before the new one: say the Telugu for to criticise, then the Telugu for to laud.
\end{quote}

\subsection*{You'll want to know: \te{వ్యాఖ్యానించు}}
\textbf{\te{వ్యాఖ్యానించు}} — \emph{vyākhyāniñcu} — \textquotedblleft{}to comment; to remark\textquotedblright{}.

\begin{quote}
\te{ఆ} \te{సినిమా} \te{గురించి} \te{అతను} \te{ఏమీ} \te{వ్యాఖ్యానించలేదు}.
\end{quote}

\begin{quote}
\emph{ā sinimā guriñci atanu ēmī vyākhyāniñcalēdu.}
\end{quote}

\begin{quote}
He made no comment about that film.
\end{quote}

Past (I …): \te{వ్యాఖ్యానించాను}, \emph{vyākhyāniñcānu}.

The noun is \te{వ్యాఖ్య}, a comment, a remark; on TV a \te{వ్యాఖ్యాత} is a commentator.

\begin{cousinweb}
\textbf{\te{వ్యాఖ్యానించు}} (\emph{vyākhyāniñcu}) is built on Sanskrit \emph{vyākhyāna}, \textquotedblleft{}explaining, exposition, gloss, comment\textquotedblright{} (Monier-Williams), with the Telugu verb ending -\te{ఇంచు}.
\end{cousinweb}

\begin{grammarlens}[title={What you've built}]
One more word for giving reasons and opinions.
\end{grammarlens}

\subsection*{Your turn}
\begin{itemize}
\raggedright
  \item \emph{Say it:} \emph{vyākhyāniñcu}
  \item \emph{Say it:} \emph{vyākhyāniñcu}, once more
  \item \emph{Say it:} \emph{ā sinimā guriñci atanu ēmī vyākhyāniñcalēdu.} — He made no comment about that film.
  \item \emph{Recall:} say the Telugu for to criticise, then the Telugu for to laud, then say \emph{vyākhyāniñcu} again
\end{itemize}

\begin{tcolorbox}[breakable,colback=teal!4,colframe=teal!35!black,title={Before you move on}]
What does \te{వ్యాఖ్యానించు} mean? (\textquotedblleft{}To comment; to remark\textquotedblright{}.) Say it once more.
\end{tcolorbox}
\section[\texorpdfstring{\te{తప్పుబట్టు} (tappubaṭṭu)}{తప్పుబట్టు (tappubaṭṭu)}]{\te{తప్పుబట్టు} (tappubaṭṭu) — to find fault with}

\label{lesson:TE-C387-tappubattu}



\begin{quote}
Before the new one: say the Telugu for to laud, then the Telugu for to comment.
\end{quote}

\subsection*{You'll want to know: \te{తప్పుబట్టు}}
\textbf{\te{తప్పుబట్టు}} — \emph{tappubaṭṭu} — \textquotedblleft{}to find fault with\textquotedblright{}.

\begin{quote}
\te{అతను} \te{ప్రతి} \te{చిన్న} \te{విషయంలో} \te{తప్పుబట్టాడు}.
\end{quote}

\begin{quote}
\emph{atanu prati cinna viṣayaṁlō tappubaṭṭāḍu.}
\end{quote}

\begin{quote}
He found fault with every little thing.
\end{quote}

Past (I …): \te{తప్పుబట్టాను}, \emph{tappubaṭṭānu}.

Two words: \te{తప్పు}, a mistake, and \te{పట్టు}, to catch. To catch someone in a mistake.

\begin{cousinweb}
\textbf{\te{తప్పుబట్టు}} (\emph{tappubaṭṭu}) is native Dravidian, a compound of \emph{tappu}, \textquotedblleft{}a mistake, a fault\textquotedblright{} (DEDR 3071), and \emph{paṭṭu}, \textquotedblleft{}to hold, to catch\textquotedblright{} (DEDR 4034), whose p sounds as b after the first word.
\end{cousinweb}

\begin{grammarlens}[title={What you've built}]
One more word for giving reasons and opinions.
\end{grammarlens}

\subsection*{Your turn}
\begin{itemize}
\raggedright
  \item \emph{Say it:} \emph{tappubaṭṭu}
  \item \emph{Say it:} \emph{tappubaṭṭu}, once more
  \item \emph{Say it:} \emph{atanu prati cinna viṣayaṁlō tappubaṭṭāḍu.} — He found fault with every little thing.
  \item \emph{Recall:} say the Telugu for to laud, then the Telugu for to comment, then say \emph{tappubaṭṭu} again
\end{itemize}

\begin{tcolorbox}[breakable,colback=teal!4,colframe=teal!35!black,title={Before you move on}]
What does \te{తప్పుబట్టు} mean? (\textquotedblleft{}To find fault with\textquotedblright{}.) Say it once more.
\end{tcolorbox}
